\documentclass[11pt,a4paper]{article}

% --- PAQUETES DE CONFIGURACIÓN ---
\usepackage[utf8]{inputenc}
\usepackage[T1]{fontenc}
\usepackage[spanish,es-tabla]{babel}
\usepackage{geometry}
\geometry{
    top=3.3cm,
    bottom=2.3cm,
    left=2.2cm,
    right=2.2cm,
    headheight=60pt,
    headsep=15pt,
    footskip=20pt
}

\usepackage{graphicx}
\usepackage{fancyhdr}
\usepackage{xcolor}
\usepackage{titlesec}
\usepackage{booktabs}
\usepackage{tabularx}
\usepackage{array}
\usepackage{amsmath}
\usepackage{enumitem}
\usepackage{microtype}
\usepackage{hyperref}

% --- PALETA INSTITUCIONAL FIEI / UNA PUNO ---
\definecolor{AzulUNA}{RGB}{10, 45, 90}
\definecolor{DoradoFIEI}{RGB}{180, 130, 20}
\definecolor{GrisOscuro}{RGB}{50, 50, 50}
\definecolor{GrisFondo}{RGB}{245, 247, 250}

% --- CONFIGURACIÓN DE LISTAS ---
\setlist{topsep=2pt, itemsep=2pt, parsep=1pt}

% --- FORMATO DE ENLACES ---
\hypersetup{
    colorlinks=true,
    linkcolor=AzulUNA,
    urlcolor=AzulUNA,
    citecolor=AzulUNA
}

% --- FORMATO DE TÍTULOS ---
\titleformat{\section}{\color{AzulUNA}\normalfont\large\bfseries}{\thesection.}{0.7em}{}[\color{AzulUNA}\titlerule]
\titleformat{\subsection}{\color{AzulUNA}\normalfont\normalsize\bfseries}{\thesubsection.}{0.5em}{}
\titleformat{\subsubsection}{\color{DoradoFIEI}\normalfont\small\bfseries}{\thesubsubsection.}{0.4em}{}
\titlespacing*{\section}{0pt}{8pt}{4pt}
\titlespacing*{\subsection}{0pt}{6pt}{3pt}
\titlespacing*{\subsubsection}{0pt}{4pt}{2pt}

% --- ENCABEZADO Y PIE DE PÁGINA (fancyhdr) ---
\pagestyle{fancy}
\fancyhf{}

% Encabezado institucional: Logo FIEI (1.45cm) y Logo Equipo (2.05cm, 40% más grande)
\fancyhead[L]{%
    \includegraphics[height=1.45cm,keepaspectratio]{../../assets/img/logo_fiei.png}%
}
\fancyhead[C]{%
    \scriptsize \textbf{\textcolor{AzulUNA}{UNIVERSIDAD NACIONAL DEL ALTIPLANO DE PUNO}}\\[2pt]
    \tiny \textcolor{GrisOscuro}{FACULTAD DE INGENIERÍA ESTADÍSTICA E INFORMÁTICA}\\[2pt]
    \scriptsize \textbf{\textcolor{AzulUNA}{Gestión de Proyectos TICS (EST315) \textbar\ Team Sobrevivientes}}\\[2pt]
    \footnotesize \textbf{\textcolor{DoradoFIEI}{ESPECIFICACIÓN FORMAL DE REQUERIMIENTOS (SRS)}}%
}
\fancyhead[R]{%
    \includegraphics[height=2.05cm,keepaspectratio]{../../assets/img/logo_equipo.png}%
}

\renewcommand{\headrulewidth}{0pt}
\renewcommand{\headrule}{}

\fancyfoot[L]{\scriptsize \textcolor{GrisOscuro}{Línea Base de Requerimientos: CasScore vs. Horus Sentinel}}
\fancyfoot[C]{\small \thepage}
\fancyfoot[R]{\scriptsize \textcolor{GrisOscuro}{Puno, Perú -- Semestre 2026-I}}
\renewcommand{\footrulewidth}{0.4pt}

\begin{document}

% =========================================================================
% PÁGINA 1: MARCO METODOLÓGICO Y REQUERIMIENTOS DE CASA SCORE
% =========================================================================

\begin{center}
    {\Large \textbf{\textcolor{AzulUNA}{ESPECIFICACIÓN FORMAL DE REQUERIMIENTOS DE SOFTWARE}}}\\[3pt]
    {\normalsize \textbf{Ingeniería de Requerimientos y Análisis Comparativo de Alternativas TIC}}\\[4pt]
    \footnotesize \textbf{Fecha de Emisión:} 04 de septiembre de 2026 \quad \textbar\quad \textbf{Versión:} 1.0 Línea Base
\end{center}

\vspace{0.1cm}

\noindent\begin{tabularx}{\textwidth}{@{}llXX@{}}
\toprule
\textbf{Curso:} EST315 -- Gestión de Proyectos TICS & \textbf{Docente:} Ing. Renzo Apaza Cutipa \\
\textbf{Equipo:} Team Sobrevivientes & \textbf{Ciclo:} Semestre Académico 2026-I \\
\bottomrule
\end{tabularx}

\vspace{0.2cm}

\section{Propósito y Marco de Ingeniería de Requisitos}
El presente documento consolida la captura y especificación técnica de requerimientos para las dos propuestas evaluadas en la Unidad I: \textbf{CASA SCORE (AlquilaPuno)} y \textbf{Horus Sentinel}. Siguiendo la premisa de \textit{Sistematizar antes de automatizar} y el ciclo IPO (\textit{Inputs $\rightarrow$ Process $\rightarrow$ Outputs}), se desglosan los requerimientos funcionales (RF), requerimientos no funcionales (RNF) y límites operativos para validar su viabilidad en un horizonte de ejecución de 6 semanas.

\section{Propuesta 1: CASA SCORE (AlquilaPuno)}

\subsection{Definición del Problema y Objetivos del Sistema}
\begin{itemize}[leftmargin=*]
    \item \textbf{Problema Central:} La migración estudiantil masiva hacia la UNA Puno genera dificultades críticas para arrendar habitaciones: comparación manual dispersa de costo, distancia y servicios, asimetría de información y vulnerabilidad ante ofertas fraudulentas o condiciones abusivas.
    \item \textbf{Objetivo General:} Desarrollar un sistema TIC inteligente de recomendación multicriterio tipo marketplace donde propietarios publiquen habitaciones y estudiantes/trabajadores encuentren opciones seguras mediante un Índice de Confianza y un sistema de reseñas bidireccional.
    \item \textbf{Objetivos Específicos:} (1) Capturar prioridades, presupuesto y cercanía al campus; (2) Normalizar y ponderar variables objetivas; (3) Generar un ranking transparente justificando el cálculo del CASA SCORE.
\end{itemize}

\subsection{Requerimientos Funcionales (RF -- CasScore)}
\noindent\begin{tabularx}{\textwidth}{p{2.3cm} p{4.6cm} X}
\toprule
\textbf{Código} & \textbf{Requisito} & \textbf{Especificación Operativa y Criterio de Validación} \\
\midrule
\textbf{RF-CS-01} & Autenticación y Perfiles & Registro y login con Supabase Auth para arrendatarios y propietarios, incluyendo perfil con validación de correo. \\
\addlinespace[3pt]
\textbf{RF-CS-02} & Publicación de Inmuebles & Registro de cuartos con variables obligatorias: canon (S/.), geolocalización, tamaño ($m^2$), servicios 24h, restricciones y fotos. \\
\addlinespace[3pt]
\textbf{RF-CS-03} & Ponderación Dinámica & Panel donde el usuario define los pesos porcentuales de sus prioridades ($w_{\text{precio}} + w_{\text{distancia}} + w_{\text{seguridad}} + w_{\text{servicios}} = 100\%$). \\
\addlinespace[3pt]
\textbf{RF-CS-04} & Normalización 0--100 & Estandarización algebraica obligatoria de variables a la escala uniforme $[0, 100]$ previa al cálculo multicriterio. \\
\addlinespace[3pt]
\textbf{RF-CS-05} & Motor de Ranking & Ejecución del algoritmo determinístico ponderado y ordenamiento interactivo de habitaciones en tiempo real. \\
\addlinespace[3pt]
\textbf{RF-CS-06} & Transparencia Algorítmica & Visualización desglosada de los puntos asignados por cada criterio (prohibición estricta de resultados en caja negra). \\
\addlinespace[3pt]
\textbf{RF-CS-07} & Índice de Confianza & Cálculo de reputación a partir de: DNI validado del propietario, fotos verificadas, transparencia de restricciones y contratos previos. \\
\addlinespace[3pt]
\textbf{RF-CS-08} & Reseñas Bidireccionales & Calificación mutua (estudiante $\leftrightarrow$ arrendador) habilitada únicamente tras confirmar el alquiler, con botones de reporte de abuso. \\
\bottomrule
\end{tabularx}

\clearpage

% =========================================================================
% PÁGINA 2: RNF CASASCORE Y REQUERIMIENTOS COMPLETOS DE HORUS SENTINEL
% =========================================================================

\subsection{Requerimientos No Funcionales y Límites (CasScore)}
\begin{itemize}[leftmargin=*]
    \item \textbf{RNF-CS-01 (Rendimiento):} El motor multicriterio debe calcular y ordenar hasta 100 alternativas habitacionales en un tiempo menor a 250 ms.
    \item \textbf{RNF-CS-02 (Stack y Despliegue Gratuito):} Frontend en Next.js (web) y prototipo móvil en Expo (React Native), backend con Supabase (PostgreSQL) y cartografía con Leaflet / OpenStreetMap en capas gratuitas.
    \item \textbf{RNF-CS-03 (Seguridad de Datos):} Implementación de Row Level Security (RLS) en PostgreSQL para aislamiento de información confidencial.
    \item \textbf{Delimitación del Alcance (No Incluye):} No realiza intermediación de pagos bancarios, no redacta contratos legales de arrendamiento, no garantiza la veracidad material de las publicaciones ni sustituye la visita física presencial.
\end{itemize}

\vspace{0.2cm}

\section{Propuesta 2: Horus Sentinel}

\subsection{Definición del Problema y Objetivos del Sistema}
\begin{itemize}[leftmargin=*]
    \item \textbf{Problema Central:} Ausencia de datos analíticos georreferenciados sobre la concentración espacial y horaria de delitos urbanos, provocando patrullajes policiales ineficientes y desplazamientos peatonales en zonas vulnerables.
    \item \textbf{Objetivo General:} Desarrollar una plataforma geoespacial predictiva que analice microdatos de denuncias e incidentes para generar mapas de calor dinámicos por zona y horario, facilitando el ruteo seguro.
    \item \textbf{Objetivos Específicos:} (1) Estimar riesgo espaciotemporal con modelos de estadística espacial; (2) Detectar brotes con concentración anómala de incidentes; (3) Desplegar rutas seguras en mapa interactivo.
\end{itemize}

\subsection{Requerimientos Funcionales (RF -- Horus Sentinel)}
\noindent\begin{tabularx}{\textwidth}{p{2.3cm} p{4.6cm} X}
\toprule
\textbf{Código} & \textbf{Requisito} & \textbf{Especificación Operativa y Criterio de Validación} \\
\midrule
\textbf{RF-HS-01} & Ingesta Espaciotemporal & Carga estructurada de denuncias: marca temporal (fecha/hora), coordenadas espaciales ($x,y$), tipo delictivo y severidad. \\
\addlinespace[3pt]
\textbf{RF-HS-02} & Densidad Kernel (KDE) & Algoritmo Kernel Density Estimation para estimar la superficie de calor continuo del riesgo delictivo basal. \\
\addlinespace[3pt]
\textbf{RF-HS-03} & Proceso de Poisson & Modelado estocástico no homogéneo para estimar réplicas delictivas condicionadas por eventos próximos en tiempo y espacio. \\
\addlinespace[3pt]
\textbf{RF-HS-04} & Capas por Franjas Horarias & Segmentación dinámica del riesgo por ventanas temporales específicas (mañana, tarde, noche y madrugada). \\
\addlinespace[3pt]
\textbf{RF-HS-05} & Detección de Anomalías & Identificación estadística de desviaciones anómalas en relación con la línea base histórica del cuadrante urbano. \\
\addlinespace[3pt]
\textbf{RF-HS-06} & Ruteo Seguro ($A^*$ Modificado) & Cálculo de trayectorias peatonales sobre grafo vial minimizando la función de coste: $\text{Costo} = \alpha \cdot \text{Distancia} + \beta \cdot \text{Riesgo}$. \\
\bottomrule
\end{tabularx}

\subsection{Requerimientos No Funcionales, Riesgos y Límites (Horus Sentinel)}
\begin{itemize}[leftmargin=*]
    \item \textbf{RNF-HS-01 (Precisión Espacial):} Georreferenciación de incidentes con tolerancia espacial máxima de 15 metros en el plano cartográfico.
    \item \textbf{RNF-HS-02 (Riesgo Clave y Mitigación de Sesgo):} Mitigar el sesgo hacia zonas históricamente más vigiladas (efecto tipo PredPol) mediante la combinación de denuncias oficiales con reportes ciudadanos auditados.
    \item \textbf{Delimitación del Alcance (No Incluye):} No tiene cobertura nacional (acotado a un distrito piloto urbano), no ejecuta reconocimiento facial, no incluye videovigilancia en vivo ni identifica sospechosos.
\end{itemize}

\clearpage

% =========================================================================
% PÁGINA 3: MATRIZ DE DECISIÓN TÉCNICA Y BLOQUE DE SUSCRIPCIÓN
% =========================================================================

\section{Matriz Comparativa y Dictamen de Factibilidad de Cartera}
Para resguardar el cumplimiento de los entregables del semestre, el equipo contrastó formalmente las especificaciones de ambas soluciones:

\vspace{0.2cm}
\noindent\begin{tabularx}{\textwidth}{p{3.8cm} X X}
\toprule
\textbf{Dimensión de Análisis} & \textbf{CASA SCORE (AlquilaPuno)} & \textbf{Horus Sentinel} \\
\midrule
\textbf{Complejidad Algorítmica} & \textbf{Media:} Normalización lineal y modelo multicriterio determinístico con suma ponderada. & \textbf{Alta:} Densidad Kernel, Poisson espaciotemporal y grafos $A^*$ con pesos multiobjetivo. \\
\addlinespace[3pt]
\textbf{Fuentes e Insumos de Datos} & \textbf{Alta e Inmediata:} Levantamiento directo de publicaciones en el entorno de la UNA Puno. & \textbf{Crítica / Bloqueante:} Requiere microdatos policiales y convenios con comisarías distritales. \\
\addlinespace[3pt]
\textbf{Factibilidad en 6 Semanas} & \textbf{Plenamente Viable:} Construcción y entrega de un MVP funcional verificado en el Sprint 5. & \textbf{Inviable:} La gestión de convenios y calibración espacial excede los plazos académicos. \\
\addlinespace[3pt]
\textbf{Costo de Infraestructura} & \textbf{Cero:} Operación completa soportada en capas gratuitas de Supabase, Vercel y Leaflet. & \textbf{Elevado:} Demanda cómputo geoespacial intensivo y servidores de procesamiento espacial. \\
\addlinespace[3pt]
\textbf{Dictamen Metodológico} & \textbf{SELECCIONADO} como proyecto oficial del equipo para el ciclo de vida de desarrollo. & \textbf{DESESTIMADO} por riesgo de bloqueo de datos externos y desborde temporal. \\
\bottomrule
\end{tabularx}

\vspace{0.3cm}

\section{Conclusión y Congelamiento de Requerimientos}
Se formaliza el congelamiento de los requerimientos de \textbf{CASA SCORE} como la Línea Base de Alcance definitiva del equipo, constituyendo el insumo obligatorio para redactar el Project Charter y la WBS.

\vspace{0.4cm}

\section{Suscripción y Conformidad del Equipo}
\noindent Los miembros del equipo técnico aprueban y suscriben la presente especificación técnica:

\vspace{1.8cm}

\begin{center}
\begin{tabular}{c c c}
    \parbox{6.5cm}{\centering \rule{6.0cm}{0.6pt}\\[5pt] \textbf{Clyde Neil Paricahua Pari}\\[2pt] {\footnotesize Project Manager}} 
    & \hspace{0.8cm} &
    \parbox{6.5cm}{\centering \rule{6.0cm}{0.6pt}\\[5pt] \textbf{Ruth Karina Apaza Solis}\\[2pt] {\footnotesize Analista de Planificación}} 
    \\[2.2cm]
    \parbox{6.5cm}{\centering \rule{6.0cm}{0.6pt}\\[5pt] \textbf{Eliseo Tarqui Ajahuana}\\[2pt] {\footnotesize Especialista Técnico / Arquitectura}} 
    & \hspace{0.8cm} &
    \parbox{6.5cm}{\centering \rule{6.0cm}{0.6pt}\\[5pt] \textbf{Viviana Mary Quisbert Quispe}\\[2pt] {\footnotesize Gestora de Comunicación y Riesgos}}
\end{tabular}
\end{center}

\end{document}